\subsection{REPL s JIT kompilací}
\label{sec:repl}

REPL je implementován v~\texttt{src/repl/} a~realizuje tzv. \emph{JIT-like}
přístup: každý vstup uživatele je přeložen do Rust kódu, zkompilován do
sdílené knihovny (\texttt{cdylib}) a~načten za běhu programu.

Vstupní řetězec prochází stejnými fázemi jako kompilace souboru:
lexing, parsing, sémantická analýza, IR lowering a~optimalizace. Následně se
použije specializovaný generátor \texttt{ReplCodegen} (\texttt{src/repl/codegen.rs}),
který vytvoří kód exportující funkci \texttt{\_\_repl\_entry}. Ta přijímá
\texttt{RuntimeContext} z~knihovny \texttt{roth-runtime} a~vykoná instrukce nad
perzistentním zásobníkem a~pamětí.

Kompilace a~dynamické načítání jsou zapouzdřeny v~\texttt{LibraryLoader}
(\texttt{src/repl/loader.rs}). Loader zapisuje vygenerovaný Rust do dočasného
adresáře, spustí \texttt{rustc --crate-type=cdylib} a~výslednou knihovnu načte
pomocí knihovny \texttt{libloading}. REPL si uchovává všechny načtené knihovny
v~paměti, aby nedošlo k~uvolnění symbolů.

Perzistentní stav mezi příkazy je rozdělen na:
\begin{itemize}
	\item \textbf{runtime kontext} (\texttt{RuntimeContext}) --- zásobník,
	      paměť proměnných a~registry slov,
	\item \textbf{kompilační kontext} (\texttt{CompilerContext}) --- slovník již
	      definovaných slov v~IR a~množinu deklarovaných proměnných. Tyto
	      informace umožňují sémantické ověření a~IR lowering s~ohledem na
	      historii REPL.
\end{itemize}
